% Generated draft: requires mathematical and source review.
\documentclass[11pt]{article}
\usepackage[margin=25mm]{geometry}
\usepackage{fontspec}
\setmainfont{Latin Modern Roman}
\usepackage{amsmath,amssymb,mathtools,unicode-math}
\setmathfont{Latin Modern Math}
\usepackage{enumitem,needspace,xurl,hyperref}
\setlength{\parindent}{0pt}
\setlength{\parskip}{5pt}
\setlength{\emergencystretch}{3em}
\newcommand{\N}{\mathbb{N}}
\newcommand{\Z}{\mathbb{Z}}
\newcommand{\Q}{\mathbb{Q}}
\newcommand{\R}{\mathbb{R}}
\newcommand{\C}{\mathbb{C}}
\newcommand{\F}{\mathbb{F}}
\newcommand{\A}{\mathbb{A}}
\newcommand{\PP}{\mathbb{P}}
\newcommand{\E}{\mathbb{E}}
\newcommand{\eps}{\varepsilon}
\newcommand{\dd}{\,\mathrm{d}}
\begin{document}
\section*{UnsolvedMath ID: 30004295}
\textbf{Source problem code:} catalogue entry
\subsection{Definitions and mathematical statement}
For an odd prime \(p\), let \(\mathcal X_p\) be the set of nonprincipal Dirichlet characters modulo \(p\). Here a Dirichlet character is the periodic extension to \(\Z\) of a homomorphism \((\Z/p\Z)^\times\to\C^\times\), with value zero on multiples of \(p\); the principal character equals one on integers coprime to \(p\). Define
\[
S_\chi(x)=\sum_{\substack{n\in\Z\\0<n\le x}}\chi(n),\qquad
M(p)=\frac{1}{\sqrt p}\min_{\chi\in\mathcal X_p}\max_{x\in[0,p]}|S_\chi(x)|.
\]
The variable \(x\) is real; the maximum is equivalently over integers \(0\le x\le p\), since these sums are constant between consecutive integers. No parity restriction is imposed on the characters.

The supplied problem asks to determine the best constant \(C\) in the strict inequality
\[
M(p)<C.
\]
The quantification over \(p\) is not explicit. Under the uniform interpretation, a precise target is to determine
\[
C_{\mathrm{unif}}=\inf\bigl\{C\in\R: M(p)<C\text{ for every odd prime }p\bigr\},
\]
with \(\inf\varnothing=+\infty\), and to specify whether this infimum itself satisfies the strict inequality for every such prime. This interpretation is conditional: the supplied wording does not distinguish a uniform bound for all primes from an eventual or asymptotic bound, or a constant depending on a fixed prime. It also does not specify whether ``best'' denotes an infimum or an attained smallest admissible constant.

There is a further boundary issue: modulo \(2\) there is no nonprincipal Dirichlet character, so the displayed minimum in the source is undefined for \(p=2\). The definitions above use odd primes to make the expression meaningful; this restriction requires confirmation against the intended source formulation.
\subsection{Short English statement}
For each odd prime, choose a nonprincipal Dirichlet character whose largest absolute initial character sum is as small as possible, and divide that maximum by \(\sqrt p\). Determine the sharp bound for this quantity. The intended quantification over primes, the meaning of a best constant for a strict inequality, and the treatment of \(p=2\) require clarification.
\subsection{Sources}
\begin{itemize}
\item Local UnsolvedMath registry: supplied record with permanent ID \(30004295\), source problem code \texttt{catalogue entry}, in the project registry context \texttt{data/unsolvedmath/problems.json} and \texttt{catalog/problems.json}. These local files were not independently inspected for this draft.
\item Supplied source citation: \emph{Analytic Number Theory (2020)}, Oberwolfach Reports, pp.~3141--3205, DOI \texttt{10.4171/owr/2019/50}. \url{https://doi.org/10.4171/owr/2019/50}. The record describes Oberwolfach Reports as documenting mathematical workshops and their open-problem discussions. It gives proposed year 2019 but supplies no individual proposer.
\item Supplied literature reference: Bober--Goldmakher--Granville--Koukoulopoulos, ``large character sums.'' \url{https://doi.org/10.4171/JEMS/799}.
\item Supplied literature reference: Bober--Goldmakher, ``distribution of maximum character sums.'' \url{https://doi.org/10.1112/S002557931200109X}.
\item Supplied literature reference: original OWR report. \url{https://doi.org/10.4171/owr/2019/50}.
\item Literature assessment dated \textbf{2026-08-22}: the supplied curation labels the problem open and states that the extremal minimum over characters in the stated P\'olya--Vinogradov normalization is not known exactly, and that distribution theorems for typical characters and general lower bounds do not determine the best constant. This is supplied dated literature triage derived from the OWR curation, not an independently verified assessment; its time-sensitive status and cited claims require independent verification.
\end{itemize}
\end{document}
